\section{Spatial Discretization}
\label{sec:spatial}

Integrating \cref{eq:ns} over a fixed control volume $\Omega_i$ and applying the
divergence theorem gives the cell-centred semi-discrete form
\begin{equation}
  V_i\frac{\mathrm{d}\overline{\mathbf{U}}_i}{\mathrm{d}t}
  = \mathbf{R}_i
  = -\sum_{f\in\partial\Omega_i}
     \left[\mathbf{H}^i(\mathbf{W}^-_f,\mathbf{W}^+_f,\mathbf{n}_f)
          -\mathbf{H}^v(\mathbf{W}_f,\nabla\mathbf{W}_f,\mathbf{n}_f)\right]\Delta s_f,
  \label{eq:fv}
\end{equation}
where $\overline{\mathbf{U}}_i$ is the cell average, $V_i$ the cell area,
$\mathbf{H}^i$ the numerical inviscid normal flux, $\mathbf{H}^v$ the viscous
normal flux of \cref{eq:viscflux}, and $\mathbf{n}_f$ the outward unit normal of
face $f$.

The implementation stores each face once and visits it once. Because the stored
normal is outward from the left cell, the accumulated flux is subtracted from the
left cell and added to the right cell:
\begin{equation}
  \mathbf{R}_{\mathrm{left}(f)} \mathrel{-}= \mathbf{H}_f\,\Delta s_f,
  \qquad
  \mathbf{R}_{\mathrm{right}(f)} \mathrel{+}= \mathbf{H}_f\,\Delta s_f .
  \label{eq:scatter}
\end{equation}
A single face therefore contributes equal and opposite amounts to its two
neighbours, which makes the discretization conservative by construction: interior
contributions telescope exactly and only boundary fluxes can change the global
integral. This is verified to $\vConservation$ in \cref{sec:verification}.

On a boundary face there is no right cell. The left state is reconstructed to the
face as usual and an appropriate exterior state is synthesised by the boundary
condition (\cref{sec:bc}); the same numerical flux function is then applied. This
keeps boundary and interior faces on one code path, so no boundary case can drift
out of sync with the interior scheme.

A face is processed if either of its cells is owned by the current rank. Faces on
a partition boundary are thus evaluated redundantly on both ranks, which is
cheaper than communicating fluxes and guarantees that the two ranks compute
bit-identical values from identical inputs.

\subsection{Second-Order Reconstruction}
\label{sec:reconstruction}

Second-order accuracy comes from a piecewise-linear reconstruction of the
\emph{primitive} variables $\mathbf{W}=(\rho,u,v,p)$. Primitives are chosen so
that the positivity of density and pressure can be enforced directly on the
reconstructed face state, which is not possible when limiting conservative
variables.

For each cell the gradient is an inverse-distance-weighted least-squares fit over
the face neighbours. The gradient $\mathbf{g}_i$ of a variable $\phi$ minimises
\begin{equation}
  J(\mathbf{g}_i)=\sum_{j\in\mathcal{N}(i)} w_j^2
  \left[\mathbf{g}_i\!\cdot\!(\mathbf{x}_j-\mathbf{x}_i)-(\phi_j-\phi_i)\right]^2,
  \qquad w_j=\frac{1}{\|\mathbf{x}_j-\mathbf{x}_i\|},
  \label{eq:lsq}
\end{equation}
whose normal equations are the $2\times2$ system
\begin{equation}
  \underbrace{\begin{bmatrix}
    \sum_j w_j^2\Delta x_j^2 & \sum_j w_j^2\Delta x_j\Delta y_j\\
    \sum_j w_j^2\Delta x_j\Delta y_j & \sum_j w_j^2\Delta y_j^2
  \end{bmatrix}}_{\mathbf{A}_i}\mathbf{g}_i
  =\sum_j w_j^2 (\phi_j-\phi_i)
  \begin{bmatrix}\Delta x_j\\ \Delta y_j\end{bmatrix} .
  \label{eq:normaleq}
\end{equation}
Since $\mathbf{A}_i$ depends only on geometry, it is inverted once during
preprocessing and folded into a per-neighbour weight vector
$\mathbf{W}_j=w_j^2\,\mathbf{A}_i^{-1}\Delta\mathbf{x}_j$, so that at run time
the gradient of every variable is the single sum
$\mathbf{g}_i=\sum_j \mathbf{W}_j(\phi_j-\phi_i)$ --- one dot product per
neighbour per variable, with no linear algebra in the hot loop. A nearly
rank-deficient stencil (possible when a cell's neighbours are almost collinear)
is detected by comparing $\det\mathbf{A}_i$ against
$10^{-12}\,\mathrm{tr}(\mathbf{A}_i)^2$ and regularised by a small diagonal
shift, rather than being allowed to produce an unbounded gradient.

Boundary faces contribute a stencil entry located at the face centroid and
carrying the \emph{physical} boundary state (\cref{sec:bc}), not the mirrored
ghost state. This distinction is essential: using a mirrored state here would
double every wall-normal gradient and corrupt the skin friction in the first cell
off the wall.

The face states are then
\begin{equation}
  \mathbf{W}^{\pm}_f=\mathbf{W}_{i}+\Phi_i\odot
  \left[\mathbf{g}_i\!\cdot\!(\mathbf{x}_f-\mathbf{x}_i)\right],
  \label{eq:recon}
\end{equation}
with $\Phi_i$ the limiter of \cref{sec:limiter} and $\odot$ a component-wise
product. Second-order reconstruction is \textbf{active in every production run
reported here}; the only first-order fallback is the positivity mechanism below,
whose activation count is recorded per run and reported in
\texttt{metadata.json}.

\subsection{Limiter and Positivity Control}
\label{sec:limiter}

Both a Barth--Jespersen and a Venkatakrishnan limiter are implemented; the
production runs use Venkatakrishnan. For each cell and variable the neighbour
envelope $\phi^{\min}_i,\phi^{\max}_i$ is gathered over the face neighbours,
including boundary-face states so that a wall-adjacent cell is not falsely
flagged as an extremum. With $\Delta_f$ the unlimited increment towards face $f$
from \cref{eq:recon} and
$y=(\phi^{\max}_i-\phi_i)/\Delta_f$ for $\Delta_f>0$ (or
$y=(\phi^{\min}_i-\phi_i)/\Delta_f$ for $\Delta_f<0$), the two limiters are
\begin{equation}
  \Phi^{\mathrm{BJ}}=\min\left(1,\,y\right),
  \qquad
  \Phi^{\mathrm{V}}=\frac{y^2+2y+\epsilon^2/\Delta_f^2}
                         {y^2+y+2+\epsilon^2/\Delta_f^2},
  \label{eq:limiter}
\end{equation}
and the cell limiter is the minimum over its faces, clipped to $[0,1]$. The
Venkatakrishnan threshold is $\epsilon^2=(K h_i)^3$ with
$h_i=\sqrt{V_i}$ the cell length scale and $K$ a case-settable constant. This
threshold is what switches the limiter off in smooth regions: without it,
limiter chattering on a stretched mesh stalls the steady residual, which is
precisely the mechanism discussed in \cref{sec:stationarity}.

Positivity is protected in two places. At the reconstruction stage, if the
limited face state still has $\rho\le0$ or $p\le0$, the whole increment is
halved repeatedly (up to eight times) and finally dropped to the first-order
cell average, which is always admissible; each such event is counted. After each
implicit update a final safety net floors density and pressure at
$10^{-8}$ of their freestream values and replaces any non-finite state by the
freestream, again counting every occurrence. Because these counts are written to
the run metadata, a run cannot quietly rely on clipping to look healthy.

\subsection{Inviscid and Viscous Fluxes}
\label{sec:fluxes}

Three inviscid numerical fluxes are implemented --- Roe with an entropy fix,
HLLC, and Rusanov/local Lax--Friedrichs --- selectable per run. Production runs
use \textbf{Roe}. With Roe averages
$\hat{\rho}=\sqrt{\rho_L\rho_R}$ and
$\hat{\phi}=(\sqrt{\rho_L}\phi_L+\sqrt{\rho_R}\phi_R)/(\sqrt{\rho_L}+\sqrt{\rho_R})$
for $\phi\in\{u,v,H\}$, and $\hat{a}^2=(\gamma-1)(\hat{H}-\tfrac12|\hat{\mathbf{u}}|^2)$,
the flux is
\begin{equation}
  \mathbf{H}^{\mathrm{Roe}}=\tfrac12\left(\mathbf{F}_L+\mathbf{F}_R\right)
  -\tfrac12\sum_{k}\tilde{\lambda}_k\alpha_k\mathbf{r}_k ,
  \label{eq:roe}
\end{equation}
with wave speeds $\lambda_{1,2,3}=\hat{u}_n-\hat{a},\ \hat{u}_n,\ \hat{u}_n+\hat{a}$,
acoustic amplitudes
$\alpha_{1,3}=(\Delta p\mp\hat{\rho}\hat{a}\Delta u_n)/(2\hat{a}^2)$, entropy
amplitude $\alpha_2=\Delta\rho-\Delta p/\hat{a}^2$, and a separately assembled
shear contribution proportional to $\hat{\rho}\,\Delta u_t$.

\paragraph{Entropy fix.} The eigenvalues are smoothed by the Harten function
\begin{equation}
  \tilde{\lambda}(\lambda,\delta)=
  \begin{cases}
    |\lambda|, & |\lambda|\ge\delta,\\[2pt]
    \dfrac{\lambda^2+\delta^2}{2\delta}, & |\lambda|<\delta,
  \end{cases}
  \label{eq:harten}
\end{equation}
applied to \emph{all three} waves. For the acoustic waves $\delta$ is the
Harten--Hyman spread, e.g.
$\delta_1=\max\!\left(0,\lambda_1-(u_{n,L}-a_L),(u_{n,R}-a_R)-\lambda_1\right)$,
which removes the non-physical expansion shock admitted by plain Roe at a sonic
point. In addition, every wave receives a floor
\begin{equation}
  \delta \ge \delta_{\min}=0.05\,
  \max\!\left(|u_{n,L}|+a_L,\ |u_{n,R}|+a_R\right).
  \label{eq:floor}
\end{equation}
The floor addresses a second, distinct failure that proved important on these
meshes. The entropy and shear waves share the eigenvalue $u_n$, which vanishes
wherever the flow is parallel to a face --- at a stagnation point, and along the
symmetry line of a symmetric body. There the density and tangential-velocity
jumps are left completely undamped, which is the classical Roe
checkerboard/carbuncle mode. It does not diverge; it quietly stalls the steady
residual and breaks the symmetry that a zero-incidence case should preserve.
Applying the floor through the same smooth Harten function keeps the flux
consistent, because identical left and right states make the jumps themselves
vanish. The coefficient $0.05$ is small enough to leave boundary layers and
contact discontinuities sharp.

\paragraph{Viscous fluxes.} Face gradients are built from the same
least-squares cell gradients. On an interior face the two cell gradients are
averaged and then corrected along the cell-to-cell direction
$\mathbf{e}=\Delta\mathbf{x}/\|\Delta\mathbf{x}\|$ so the face gradient
reproduces the actual cell-value difference exactly,
\begin{equation}
  \nabla\phi_f=\overline{\nabla\phi}
  +\left(\frac{\phi_R-\phi_L}{\|\Delta\mathbf{x}\|}
  -\overline{\nabla\phi}\cdot\mathbf{e}\right)\mathbf{e},
  \label{eq:facegrad}
\end{equation}
which suppresses the odd-even decoupling that a plain average allows. The
temperature gradient is derived from the density and pressure gradients through
$T=p/(\rho R)$,
$\nabla T=R^{-1}\left(\nabla p/\rho-p\,\nabla\rho/\rho^2\right)$.

On a wall face the viscous flux is evaluated with the \emph{physical} wall state
and the wall-normal derivative is corrected using the imposed wall value: with
$d_n=(\mathbf{x}_f-\mathbf{x}_i)\cdot\mathbf{n}$,
\begin{equation}
  \nabla\phi \leftarrow \nabla\phi
  +\left(\frac{\phi_{\mathrm{wall}}-\phi_i}{d_n}
  -\nabla\phi\cdot\mathbf{n}\right)\mathbf{n}.
  \label{eq:wallgrad}
\end{equation}
This is what makes the skin friction reflect the actual near-wall profile on a
stretched boundary-layer mesh instead of the much smaller cell-average gradient.
For an adiabatic wall the normal component of $\nabla T$ is then projected out
exactly, enforcing $\partial T/\partial n=0$.

\paragraph{Positivity fallback.} If the Roe average is inadmissible
($\hat{a}^2\le0$, possible for a very strong expansion), the flux falls back to
HLLC for that face, which is positivity preserving. This is reported in
\texttt{metadata.json} as part of the flux description rather than hidden. In
practice this path is never taken in any production run reported here: every case
records zero positivity-fallback events. That has a consequence for what the
production results can and cannot verify, discussed in \cref{sec:hllcbug}.

\paragraph{Named dissipation and limiter parameters.} For completeness, the
scheme carries four tuned constants, all disclosed here rather than buried in the
source: the linear-wave dissipation floor coefficient $\pLinearFloor$ in
\cref{eq:floor}; the Venkatakrishnan threshold constant $K=\pVenkatK$ in
\cref{eq:limiter}; and the viscous spectral-radius factors $\pViscTimeFactor$ in
the local time step \cref{eq:dtau} and $\pViscJacFactor$ in the implicit
off-diagonal \cref{eq:offdiag}.

\subsubsection{Added dissipation on the linear waves: disclosure}
\label{sec:lineardissipation}

The floor \cref{eq:floor} deserves to be stated plainly, because only part of it
is textbook. Smoothing the \emph{acoustic} eigenvalues by the Harten--Hyman
spread is the standard entropy fix and needs no defence. Applying a floor to the
\emph{entropy and shear} eigenvalues is different: it is added artificial
dissipation, a carbuncle cure, and it is not part of the classical Roe scheme.

It is present because without it these cases do not converge correctly. The
entropy and shear waves share the eigenvalue $u_n$, which vanishes identically on
any face where the flow is parallel to the face --- including the stagnation
streamline of a symmetric body at zero incidence, which is exactly the
configuration of all six airfoil cases. There, density and tangential-velocity
jumps are advected with zero dissipation, producing a checkerboard mode that does
not diverge but stalls the residual and destroys the symmetry the solution should
have.

The cost is quantifiable. At a stagnation face where the exact numerical flux is
$[0,\,p,\,0,\,0]$ in normalised form, the floored scheme returns
$[\num{-2.96e-4},\,1.0,\,\num{2.13e-4},\,\num{7.42e-5}]$: a mass-flux error of
about $3\times10^{-4}$ relative to the pressure term. The scheme remains
consistent, since identical left and right states make all jumps vanish and the
exact physical flux is recovered, but at a stagnation point the added dissipation
is genuinely there.

The honest concern is that at $\Reinf=5000$ this numerical dissipation competes
with the physical viscosity, so a fraction of the computed skin friction could in
principle be an artefact of the floor rather than of $\mu$. That concern was tested
rather than argued away, by sweeping the coefficient in controlled experiments.
\Cref{sec:floorsens} reports the outcome in full, and it does not support the
design choice: the floor changes $C_D$ by $0.021\,\%$ on the cylinder, and
disabling it entirely produces no checkerboard mode even on the zero-incidence
inviscid airfoil, where the entropy and shear eigenvalues vanish on the stagnation
streamline and no physical viscosity is present to damp a carbuncle. The floor is
therefore retained as a safeguard for conditions not exercised here, and no result
in this report depends on it.
